% ---------------------------------------------------------------------------
% FRNN: the original grid algorithm and the libFRNN redesign
%
% Self-contained and compilable, but every section is also safe to \input into
% a larger document: the preamble uses only amsmath, hyperref and url.
% ---------------------------------------------------------------------------
\documentclass[11pt]{article}

\usepackage[margin=1in]{geometry}
\usepackage{amsmath}
\usepackage{amssymb}
\usepackage{url}
\usepackage[hidelinks]{hyperref}

\title{Fixed-Radius Nearest Neighbors on the GPU:\\
       From the Original FRNN Grid to \texttt{libFRNN}}
\date{}

\begin{document}
\maketitle

% ===========================================================================
\section{The Original FRNN Algorithm}
\label{sec:original}

% ---------------------------------------------------------------------------
\subsection{Original Use Case}
\label{sec:original-usecase}

The original FRNN implementation\footnote{\url{https://github.com/lxxue/FRNN}}
was written as a drop-in accelerator for point-cloud deep learning, exposing an
interface deliberately modeled on \texttt{pytorch3d.ops.knn\_points}. The target
workload is a batch of variable-length 3D point clouds --- registration, scene
flow, correspondence matching, and implicit-surface reconstruction --- in which
a network repeatedly asks, for every point of a query cloud $P_1$, for the $K$
nearest points of a reference cloud $P_2$ that lie within a fixed radius $r$.
Two properties of that setting shape the entire design. First, the radius is
fixed and known before the search, which makes a uniform spatial grid the
natural acceleration structure. Second, the reference cloud is frequently
static across many calls (the same scene queried by successive network layers
or successive iterations), which is why the library returns its grid as a
cacheable handle and can skip reconstruction entirely on subsequent calls. The
implementation was written for $D \in \{2,3\}$, differentiable with respect to
the input coordinates, and everything above the CUDA kernels --- grid sizing,
buffer allocation, batching --- lives in Python and PyTorch tensors.

% ---------------------------------------------------------------------------
\subsection{Method}
\label{sec:original-method}

\paragraph{The uniform grid as a filter.}
The algorithm replaces the $O(|P_1| \cdot |P_2|)$ all-pairs comparison with a
two-stage \emph{bin-and-scan} scheme. A uniform grid is laid over the axis-aligned
bounding box of the reference cloud with a cell edge of $r/\rho$, where $\rho$
is the \texttt{radius\_cell\_ratio} (default $2$); the resolution per axis is
clamped to $64$ in 3D and $512$ in 2D, and the cell size is enlarged if that cap
would be exceeded. Because a cell edge is at most $r$, every point within
distance $r$ of a query $q$ must lie in the cell block
\[
  \Big[\big\lfloor (q_d - m_d - r)\,\delta \big\rfloor,\;
       \big\lfloor (q_d - m_d + r)\,\delta \big\rfloor \Big],
  \qquad d = 0,\dots,D_{\text{grid}}-1,
\]
with $m$ the grid minimum and $\delta = 1/\text{cell}$ the inverse cell size.
The grid is therefore a \emph{conservative filter}: it never discards a true
neighbor, and the exact test $\lVert p - q \rVert^2 \le r^2$ is still applied to
every candidate the stencil produces. Grid parameters are packed into an
$8$-float record in 3D (minimum, $\delta$, per-axis resolution, total cells) and
a $6$-float record in 2D.

\paragraph{Building the grid without a sort.}
Construction is three kernels and no comparison sort. \texttt{InsertPoints}
maps each point to a linear cell index
$c = (g_x R_y + g_y) R_z + g_z$ and performs an \texttt{atomicAdd} on that
cell's counter; the value \emph{returned} by the atomic is kept as the point's
rank within its cell, so one pass simultaneously builds the cell histogram and
assigns each point a unique, stable slot. An exclusive prefix sum over the cell
counts --- a work-efficient Blelloch scan with the classical
\texttt{CONFLICT\_FREE\_OFFSET} shared-memory padding to avoid bank conflicts
--- converts counts into per-cell base offsets. \texttt{CountingSort} then
writes each point to position $\text{offset}[c] + \text{rank}$, producing a
coordinate array in which the membership of any cell is a \emph{contiguous}
range $[\text{off}[c], \text{off}[c{+}1])$, plus a permutation array mapping
sorted positions back to original indices. The candidate stream for a stencil
cell is thus a single linear read, which is exactly what a GPU wants.

\paragraph{Sorting the queries too.}
A subtle but important choice: the \emph{query} cloud is inserted, prefix-summed
and counting-sorted through the same three kernels. Queries are then processed
in grid order, so the threads of a warp are spatially adjacent and traverse
overlapping stencils; the candidate ranges they stream are largely the same
cache lines. Results are scattered back to the caller's original indexing
through the query permutation, so the reordering is invisible at the API.

\paragraph{Bounded selection in registers.}
Each thread owns one query and maintains a \texttt{MinK} structure --- a
bounded array of the $K$ smallest squared distances seen so far, together with
their indices, plus a cached maximum. Insertion is $O(1)$ when the candidate is
rejected or the structure is not yet full, and $O(K)$ when it displaces the
current maximum (the maximum is re-found by a linear rescan). The structure
allocates nothing: it is constructed over arrays supplied by the caller, which
are thread-local registers in the fast kernels and a slice of the output tensor
in the generic one. A final in-place sort orders each row by distance.

\paragraph{Grids over a coordinate prefix.}
The grid is never built over all $D$ axes. In two dimensions it is a 2D grid; in
every dimension $D \ge 3$ it is a \emph{3D} grid over the first three
coordinates, with the remaining $D-3$ coordinates participating only in the
exact distance evaluation inside the innermost loop. (Later work in this line
widens that prefix to four axes for high-dimensional inputs; see
Section~\ref{sec:libfrnn}.) The correctness argument is that projection onto a
coordinate subset can only shrink distances,
\[
  \lVert p - q \rVert_{\text{proj}} \;\le\; \lVert p - q \rVert,
\]
so a stencil computed in the projected space is still a superset of the true
neighborhood; only its \emph{selectivity} degrades as $D$ grows. This is what
keeps the cell count bounded --- a full $D$-dimensional grid would need
$\text{res}^{D}$ cells and a $3^{D}$ stencil --- at the price of an increasingly
weak filter in high dimensions.

\paragraph{CUDA-level memory management.}
Several deliberate choices target the memory system rather than arithmetic.
(i)~Every kernel uses a fixed launch geometry ($256$ blocks $\times$ $256$
threads) with a chunked persistent loop over the work items, so occupancy is
decoupled from $N$ and no launch is oversized for a small cloud. (ii)~All
pointer arguments are marked \texttt{\_\_restrict\_\_}, allowing the compiler to
route the read-only point and offset arrays through the texture/constant cache
path. (iii)~The $D=2$ and $D=3$ specializations load coordinates as
\texttt{float2}/\texttt{float3} vector types, turning per-axis scalar loads into
single 8- or 12-byte transactions. (iv)~Kernels are templated on $D$ and $K$
through a recursive dispatch ladder and instantiated in three tiers: a fully
register-resident variant for small $D$ and $K$, a middle variant that caches
the query point in registers but writes selections straight into the output
tensor, and a generic variant for large $D$; compile-time $D$ and $K$ let the
inner loops unroll completely and remove all indirect indexing. (v)~The
selection structure writes directly into the destination tensor in the generic
tiers, avoiding a second pass over results. (vi)~Above the kernels, the built
grid is returned to the caller and can be handed back on the next call, so a
static reference cloud pays construction cost exactly once. A brute-force path
is retained for clouds small enough that grid construction does not amortize.

% ---------------------------------------------------------------------------
\subsection{Possible Improvements}
\label{sec:original-improvements}

\paragraph{Entanglement with PyTorch.}
The most consequential limitation is architectural rather than numerical: the
algorithm is not separable from its framework. Grid sizing runs in a Python loop
over the batch, calling \texttt{.min()}, \texttt{.max()} and \texttt{.item()} on
device tensors --- each of which is a host--device synchronization --- and then
copies the parameter tensor back to the host with \texttt{.cpu()} before the
prefix sum, which is itself launched once per batch element from Python. A
search that should be a handful of asynchronous launches instead becomes a
sequence of stalls whose cost is invisible in kernel profiles and dominant in
wall-clock latency at small $N$. Everything crosses the ATen dispatcher, the
extension is ABI-pinned to a specific PyTorch build, and the kernels are
unreachable from a C++ simulation, a NumPy program, or a CUDA graph.

\paragraph{Allocation churn and workspace reuse.}
Every call allocates a fresh set of scratch buffers: per-cell counts and
offsets, per-point cell and rank arrays, the sorted coordinate array, and the
sorted-index permutation --- for both clouds. The count and offset arrays are
sized by the \emph{total} cell count, so a $64^3$ grid costs $262{,}144$
integers per batch element regardless of how few cells are occupied. Nothing is
retained between calls, and the query-side grid is rebuilt even when only the
query moved. A persistent, monotonically growing workspace would remove the
allocator from the timed path entirely. Similarly, the hand-rolled Blelloch scan
predates the availability of a tuned device-wide scan and could be replaced by a
single vendor-optimized call.

\paragraph{Memory layout at high $D$.}
Both point arrays are array-of-structures, indexed \texttt{p[i*D + d]}. A warp
reading one axis therefore touches $32$ distinct cache lines --- at $D=16$ the
inner loop fetches roughly an order of magnitude more bytes than it consumes.
A structure-of-arrays layout for the sorted reference cloud, \texttt{p[d*P + i]},
packs $32$ candidates per cache line for a given axis and is the single highest
leverage change for $D \ge 8$.

\paragraph{Selection cost and grid selectivity.}
\texttt{MinK}'s displacement path rescans all $K$ entries to re-find the
maximum, so dense queries at $K = 64$ or $128$ spend most of their time in
selection rather than in distance evaluation; a heap has the right asymptotics
but the wrong constant for sparse queries, which suggests a hybrid. Independently,
the fixed three-axis projection becomes a weak filter as $D$ grows: cells fill
with candidates that pass the projected test and fail the full-dimensional one,
and no cell-level bounding-box rejection exists to discard stencil corners that
cannot intersect the query ball.

\paragraph{Register pressure and compile time.}
Selection buffers are declared at the template's maximum $K$ even when the call
uses a fraction of it, and the query cache occupies a further $D$ registers, so
the upper end of the ladder spills. Because the ladder instantiates the full
$D \times K$ cross product, build times are long enough that the project's own
installation notes warn about them.

\paragraph{Synchronous, stream-oblivious API.}
There is no way to enqueue a search on a caller-supplied stream, to run it
count-first and allocate an exactly sized output, or to keep the result on the
device. The output itself is a dense $[N, K]$ tensor padded with sentinels,
which the caller must post-process to obtain the edge list that downstream graph
code actually wants.

% ===========================================================================
\section{\texttt{libFRNN}: A Standalone Redesign}
\label{sec:libfrnn}

The improvements enumerated in Section~\ref{sec:original-improvements} are, one
for one, the design brief of
\texttt{libFRNN}\footnote{\url{https://github.com/xju2/libFRNN}}. It keeps the
part of the original that is genuinely good --- bin by \texttt{atomicAdd}, prefix
sum, counting sort, scan a conservative stencil, select in registers --- and
rebuilds everything around it: the core becomes a standalone C++17/CUDA library
with no dependency on Python, PyTorch, ATen, \texttt{c10}, TBB or NumPy; grid
sizing moves onto the device so no call synchronizes with the host; all scratch
memory moves into a reusable workspace so the warm path performs no allocation;
the reference cloud is stored structure-of-arrays above four dimensions; the
grid is widened from three to four indexed axes with per-cell bounding-box
rejection to pay for it; selection switches adaptively between an insertion list
and a deterministic heap; and the whole search is expressed as stream-ordered
asynchronous work with a count-first/materialize split, under an explicit
exactness and tie-breaking contract that makes the aggressive optimizations
provably safe. The remainder of this section takes those changes in turn.

\paragraph{A core that owes nothing to a framework.}
The public surface is a handful of plain structures --- \texttt{PointView},
\texttt{DevicePointView}, \texttt{Edge}, \texttt{BuildOptions} --- and free
functions in namespace \texttt{frnn}, installed as a CMake package consumed with
\texttt{find\_package(frnn)}. The NumPy binding is a thin pybind11 translation
layer, and PyTorch support is a \emph{separate, optional} build target that
registers a custom operator, reads \texttt{data\_ptr}, enqueues on PyTorch's
current stream, and returns a CUDA tensor without ever touching host memory.
Importing the base package does not import PyTorch; the same object code serves
a C++ reconstruction pipeline, a NumPy script, and a training loop. In practice
this also removes the ABI pin: the library is rebuilt against a CUDA toolkit,
not against a specific tensor framework release.

\paragraph{Grid construction on the device.}
The Python loop that computed bounds and cell sizes is replaced by three
kernels. \texttt{initializeBounds} seeds a bounding-box record;
\texttt{computeBounds} reduces the reference cloud into it using
\texttt{atomicMinimum}/\texttt{atomicMaximum} built from compare-and-swap on the
bit pattern of the float; \texttt{finalizeGrid} runs on a single thread and
derives the cell size and per-axis resolution, iterating up to eight times to
coarsen the cell size until the cell budget is respected. The grid parameters
never leave the device, so a search is a pure sequence of enqueues: no
\texttt{.item()}, no \texttt{.cpu()}, no implicit synchronization anywhere in the
build phase.

\paragraph{A persistent workspace, and scans that are not hand-written.}
Every scratch buffer --- bounds, grid parameters, cell counts, cell offsets,
per-point cells and ranks, the sorted cloud, the permutation, neighbor distances
and indices, edge counts and offsets, and the scan temporary --- lives in a
\texttt{Workspace} whose \texttt{reserve} grows capacities monotonically and
returns immediately when the existing allocation suffices. After the first call
of a given shape the warm path issues \emph{zero} \texttt{cudaMalloc} calls, and
clearing is device-guarded so a workspace frees on the device that owns it. Both
prefix sums --- over cell counts and over per-query edge counts --- are
\texttt{cub::DeviceScan::ExclusiveSum}, sized once and sharing a single temporary
buffer allocated at the maximum of the two requirements. Every memset is
\texttt{cudaMemsetAsync} on the caller's stream.

\paragraph{A four-axis grid, paid for by corner rejection.}
The indexed prefix is now dimension-dependent: $\min(D,3)$ axes for $D \le 4$ and
\emph{four} axes for $5 \le D \le 32$, with separate cell budgets ($128{,}000$
cells for the 3D grid, $128^3$ for the 4D grid) and a cell edge of $r$ in three
dimensions but $r/2$ in four. Widening the projection strengthens the filter
exactly where the original was weakest, but it also turns a $3^3$ stencil into a
$3^4$ one, most of whose cells are corners that cannot intersect the query ball.
The search kernel therefore computes, for four-dimensional grids only, the
per-axis gap between the query and the cell's bounding box, accumulates the
squared gap axis by axis, and \texttt{continue}s as soon as it exceeds $r^2$ ---
so a corner cell is rejected after one or two multiplies instead of a full scan
of its contents. The same test is deliberately \emph{not} applied to the 3D
stencil, where its arithmetic would cost more than it saves.

\paragraph{Structure-of-arrays storage and screened distance evaluation.}
\texttt{countingSort} writes the sorted reference cloud AoS through dimension
four --- where a candidate record fits in one short contiguous burst --- and SoA
above it, at offset $d \cdot P + i$. The high-dimensional distance routine
exploits that layout twice. It first performs a \emph{screen}: it accumulates
the squared difference over groups of four axes, starting from axis four and
wrapping around, and bails out as soon as the partial sum exceeds an
\emph{expanded} threshold. Candidates that survive the screen are then
recomputed in canonical axis order with fused multiply-adds. Because the screen
threshold is widened beyond the worst-case rounding difference between two
non-negative float32 sums of at most $32$ terms, a candidate near the boundary is
always recomputed canonically, and the reordering therefore cannot change either
radius membership or the top-$K$ ordering. The performance argument is a cache
argument: for a given axis, SoA packs $32$ candidates per cache line against
roughly $2.67$ for AoS, which is what makes an early-exit filter worth running
at all.

\paragraph{Hybrid selection with a total order on ties.}
Selection begins as the original's sorted insertion list, which is optimal for
the sparse queries that dominate a well-chosen radius. Once a query has accepted
$24$ candidates and $K \ge 64$, the list is heapified in place and the query
switches to a deterministic max-heap, with a final heap-sort to restore sorted
order --- so dense queries stop paying the quadratic shifting cost and sparse
queries never pay heap overhead. Both paths order candidates by the strict total
order $(\text{squared distance}, \text{original index})$, which removes the
original's unspecified behavior on equal distances: ties resolve to the lower
index, always, on every architecture and along either selection path.

\paragraph{Runtime dispatch instead of a template cross-product.}
Dispatch is a \texttt{switch} on the runtime dimension over the specialized set
$\{1,2,3,4,8,12,16\}$ with a generic compiled path covering the remainder through
$D = 32$; $K$ stays a runtime value. This collapses the original's $D \times K$
instantiation matrix to a short list, cutting compile time drastically while
retaining full unrolling where it matters. Occupancy is managed explicitly:
$256$-thread blocks for $D \le 4$, $128$-thread blocks with
\texttt{\_\_launch\_bounds\_\_(128, 8)} for the high-dimensional kernels, which
caps register allocation at a level that keeps eight blocks resident per SM. The
original's query-reordering trick is preserved --- when query and reference are
the same set, queries are visited in the sorted spatial order --- and results are
still written back by original index. Algorithm choice itself is automatic,
through an overflow-safe work estimate
$Q \cdot P \cdot D \cdot \lceil K/8 \rceil \le 2\times10^6$ that routes small
problems to exact brute force, with a manual override available for measurement.

\paragraph{An asynchronous, count-first device API.}
\texttt{buildEdgesAsync} takes device pointers, a caller-owned output buffer, a
workspace and a stream, and enqueues every memset, scan and kernel on that
stream without synchronizing it. Passing a null edge pointer requests
\emph{count-only} execution: the exact device-resident edge count is produced,
and after the caller synchronizes on it, \texttt{materializeEdgesAsync} writes
the already-prepared edges into an exactly sized allocation without repeating
the search. Launch and API errors are reported immediately; execution errors
surface at the caller's synchronization point. The synchronous host API is built
on top of the same path and is the only place validation of host-side
coordinates occurs, precisely because inspecting device values on the host would
destroy the asynchrony.

\paragraph{Edges as the output, and an explicit exactness contract.}
Rather than returning a padded $[N,K]$ neighbor tensor for the caller to
post-process, the library counts per-query edges, exclusive-scans the counts, and
writes a compacted $[E,2]$ edge array in one pass --- with self-loop removal and,
for the single-set case, one emission per undirected pair under the convention
$\text{source} > \text{target}$. Underpinning all of it is a written contract:
squared distances accumulate in float32 in dimension order with one fused
multiply-add per coordinate; the radius boundary is inclusive; rows are ordered
by $(\text{distance}, \text{index})$ and truncated to $K$; cell ranges, cell
boxes and rejection thresholds are conservatively expanded to absorb float32
roundoff so that no true neighbor can be lost to a boundary computation. That
contract is validated against an independent CPU oracle using \texttt{std::fma},
with exact-boundary, adjacent-\texttt{nextafter} and rounding-sensitive
high-dimensional cases. It is what licenses the screening, the early exits and
the bounding-box rejection: each is a pure performance transformation that
provably cannot alter the result.

\end{document}
